\section{Analysis and Limitations}

\paragraph{What the gate labels mean.}
V-SUP means that a raw baseline has a supported positive pressure-specific gap under the \auditname matched-neutral contract. DynaGuard and WildGuard meet this criterion and are the only selected-main baselines used as attenuation evidence. V-NS means that the current 50-base protocol does not support a positive raw vulnerability. BingoGuard and HarmAug are therefore not treated as failed baselines or as removal evidence; instead, they act as guardrail baselines where the claim gate asks whether \method preserves their status and avoids introducing a new supported residual gap. R-NS means that the post-projection residual pressure gap is not supported.

\paragraph{Why the projection should not be oversold.}
\method replaces the pressure-side score with a matched neutral-control estimate for audit evaluation. It asks whether the supported raw gap is accounted for by the matched-neutral readout under the archived artifact contract. This is the right operation for a residual-gap audit, but it is not the same as observing a one-pass model ignore pressure at deployment time. It does not establish original-judge pressure invariance, production deployment readiness, equal-cost mitigation, or unique estimator optimality. Any use outside audit settings would need an explicit neutral-control generation or cache policy and a cost accounting that is not provided here.

\paragraph{What the selective wrapper adds.}
The neutral-consensus selective wrapper is the only experiment in the paper that behaves like an intervention rather than a pure residual-gap readout: it either issues a neutral-consensus automatic decision or abstains/escalates. This helps answer the practical objection that the audit is merely post-hoc. However, it is still an offline replay over archived artifacts. It requires extra matched-control queries, assumes authored neutral controls are suitable at test time, and depends on the same construction validity that has not yet received blinded human validation. We therefore report coverage, abstention, correction, and induced-error rates rather than accuracy gains, and we do not treat the wrapper as deployable, equal-cost, single-pass, or a replacement for \method's claim gate.

\paragraph{Neutral-control validity.}
The neutral controls are authored to preserve the same base case, label, layout, and format while removing the pressure cue, and the artifact checks coverage, field-structure matching, and token-difference summaries. We now add an aggregate neutral-control validity audit over the 50-base main gate, PKU200 scale audit, and non-PKU 200-base source-pair audit: across 450 bases, 19,350 records, 1,800 attack cells, and 4,500 matched neutral records, the audit finds zero missing matched-neutral cells, zero neutral-role failures, zero metadata mismatches, and zero scanner hits for explicit desired-label or social-pressure cue leakage. Clean-versus-neutral prediction sanity checks show small neutral-minus-clean error shifts across the archived DynaGuard and WildGuard runs. These are authored construction checks, mechanical artifact checks, and behavioral sanity checks, not an independent validation study. We still do not have blinded human judgments showing that every neutral control preserves perceived safety label and difficulty while removing the pressure cue. A future audit should sample matched cells, hide role assignment from annotators, and ask whether label preservation, difficulty preservation, and pressure removal hold with inter-annotator agreement. Without that study, the paper should claim a matched-neutral audit contract with mechanical validity evidence, not validated counterfactual equivalence.

\paragraph{What the ablations imply.}
The no-inference projection ablations make the method claim narrower, not broader. They rule out arbitrary neutral substitution: global neutral and same-cell other-base projections create large supported residual gaps for DynaGuard and WildGuard. But they also show that same-base controls that ignore exact layout or format can attenuate the gap as much as, or more than, matched mean-v1 on this split. This means the current evidence supports a base-matched audit framework and fail-closed inference contract more strongly than it supports a unique estimator choice. Future work should preregister estimator variants and evaluate them on fresh sources before claiming an optimal projection rule.

\paragraph{WildGuard sparsity.}
The WildGuard attenuation result is real but thin. The paired primary mean attenuation is positive and Holm-significant, but the base-level distribution is sparse: 7 positive bases, 42 zero bases, and 1 negative base for the primary gap, with median 0. The clean-correct flip attenuation is also sparse: 6 positive bases and 36 zero bases. The paper therefore claims positive paired mean attenuation, not a broad per-base dominance result.

\paragraph{Selected-baseline scope.}
The current main claim uses DynaGuard, BingoGuard, HarmAug, and WildGuard because each is in the frozen accepted-baseline policy and has a complete raw and post-hoc artifact under the same 50-base contract. ShieldLM is diagnostic only because its official response-safety contract does not directly match the rendered safety-judge review setting used by \auditname; its supplementary result is disclosed but cannot support or refute the selected-main claim. BeaverDam/R2/ShieldGemma are not part of the current complete evidence ledger. This prevents cherry-picking from being hidden, but it also limits external validity.

\paragraph{Threshold and slice robustness.}
The formal slice-inference report shows that DynaGuard's raw effect appears in multiple family and layout slices, while WildGuard's raw hard-label effect is concentrated in the toward-unsafe label/direction stratum and does not survive family- or layout-level Holm correction. DynaGuard and WildGuard expose mostly hard-label scores, so threshold checks at 0.25, 0.50, and 0.75 are not a strong continuous-score robustness test; they mainly confirm that the reported 0.5 threshold is not hiding a soft-score tuning choice for these runs. HarmAug supplies the current continuous-score diagnostic: it shows small within-field Holm-supported adverse-probability drift in several raw slices, while its hard-label gate remains V-NS and the projected adverse-probability gaps collapse to zero. The regenerated artifact now includes run-metric, metric-global, and report-global Holm screens in addition to local within-field Holm. These screens reduce the risk of cherry-picked slice language, but they still do not turn slice localization into a new primary claim gate.

\paragraph{Dataset scale and pressure families.}
The 50-base split is intentionally controlled and fully paired, but it is underpowered for broad claims and its selected-main claim gate is tied to the frozen 50-base contract. The paper treats this gate as a controlled audit, not as evidence for source-general, model-general, or full-dataset robustness. The 200-base PKU expansion and PKU2K full-data diagnostic strengthen raw-vulnerability replication for DynaGuard and WildGuard, but they also narrow the method claim: mean-v1 attenuates both baselines substantially while leaving small statistically supported residual gaps. This is the intended role of the scale diagnostics. They reduce the risk that the phenomenon is a 50-base artifact, while the fail-closed gate prevents those same diagnostics from being converted into a broader residual-elimination or full-data main claim. The non-PKU 200-base source-pair audit replicates raw vulnerabilities and strong paired attenuation, but source and label are coupled, so it cannot support general source robustness or safe/unsafe asymmetry. The PKU2K cycle-v1 diagnostic has near-zero residual gaps, but it was introduced after observing the mean-v1 residual and is therefore a candidate for future preregistration rather than a selected-main result. Future work should add source pairs with both labels per source, a main-compatible continuous-score baseline, preregistered estimator variants, and additional pressure styles, languages, and guard models that were not available in the present artifact ledger.

\paragraph{What text cannot fix.}
Several reviewer concerns are evidence limitations rather than wording problems. The frozen 50-base gate remains the primary sample size and is deliberately controlled rather than broad. The 200-base and PKU2K diagnostics support raw-vulnerability replication and attenuation, but not universal residual elimination or a replacement primary result. The non-PKU source pair cannot establish source-general robustness because source and label are coupled. The hard-label adapters cannot support strong calibration or probability-threshold claims; they mainly support label-level error-gap analysis, with HarmAug as the only current continuous-score diagnostic. Finally, no textual revision, and no offline selective-wrapper replay, can turn a post-hoc matched-control audit into an equal-cost, single-pass, deployable guard.

\paragraph{Claim boundary.}
The supported claim is the selected-main-baseline \auditname post-hoc audit claim. It does not support unrestricted superiority across safety-judge settings, single-pass robustness, trained PACT, equal-cost comparison against one-pass guards, or generalization beyond the matched-neutral contract.
